% Options for packages loaded elsewhere
\PassOptionsToPackage{unicode}{hyperref}
\PassOptionsToPackage{hyphens}{url}
\documentclass[
  12pt,
]{ctexart}
\usepackage{xcolor}
\usepackage{amsmath,amssymb}
\setcounter{secnumdepth}{-\maxdimen} % remove section numbering
\usepackage{iftex}
\ifPDFTeX
  \usepackage[T1]{fontenc}
  \usepackage[utf8]{inputenc}
  \usepackage{textcomp} % provide euro and other symbols
\else % if luatex or xetex
  \usepackage{unicode-math} % this also loads fontspec
  \defaultfontfeatures{Scale=MatchLowercase}
  \defaultfontfeatures[\rmfamily]{Ligatures=TeX,Scale=1}
\fi
\usepackage{lmodern}
\ifPDFTeX\else
  % xetex/luatex font selection
\fi
% Use upquote if available, for straight quotes in verbatim environments
\IfFileExists{upquote.sty}{\usepackage{upquote}}{}
\IfFileExists{microtype.sty}{% use microtype if available
  \usepackage[]{microtype}
  \UseMicrotypeSet[protrusion]{basicmath} % disable protrusion for tt fonts
}{}
\makeatletter
\@ifundefined{KOMAClassName}{% if non-KOMA class
  \IfFileExists{parskip.sty}{%
    \usepackage{parskip}
  }{% else
    \setlength{\parindent}{0pt}
    \setlength{\parskip}{6pt plus 2pt minus 1pt}}
}{% if KOMA class
  \KOMAoptions{parskip=half}}
\makeatother
\usepackage{longtable,booktabs,array}
\usepackage{calc} % for calculating minipage widths
% Correct order of tables after \paragraph or \subparagraph
\usepackage{etoolbox}
\makeatletter
\patchcmd\longtable{\par}{\if@noskipsec\mbox{}\fi\par}{}{}
\makeatother
% Allow footnotes in longtable head/foot
\IfFileExists{footnotehyper.sty}{\usepackage{footnotehyper}}{\usepackage{footnote}}
\makesavenoteenv{longtable}
\setlength{\emergencystretch}{3em} % prevent overfull lines
\providecommand{\tightlist}{%
  \setlength{\itemsep}{0pt}\setlength{\parskip}{0pt}}
\usepackage{bookmark}
\IfFileExists{xurl.sty}{\usepackage{xurl}}{} % add URL line breaks if available
\urlstyle{same}
\hypersetup{
  hidelinks,
  pdfcreator={LaTeX via pandoc}}

\author{SCX}
\date{2026}

\begin{document}

\section{Online Learning Regret Analysis of SCX
Gatekeeper}\label{online-learning-regret-analysis-of-scx-gatekeeper}

\begin{quote}
\textbf{Version}: 2026-06-28 \textbar{} \textbf{Status}: Theoretical
development \textbar{} \textbf{Audit}: Pre-verification
\textbf{Purpose}: Analyze the SCX gatekeeper's online learning regret
under delayed feedback from the NEP student, derive regret bounds using
standard online convex optimization tools, and discuss the impact of the
growing memory bank. \textbf{Prerequisites}: Document 01 (Symbol System
and Problem Setup)
\end{quote}

\begin{center}\rule{0.5\linewidth}{0.5pt}\end{center}

\subsection{Table of Contents}\label{table-of-contents}

\begin{enumerate}
\def\labelenumi{\arabic{enumi}.}
\tightlist
\item
  \hyperref[1-problem-formulation]{Problem Formulation}
\item
  \hyperref[2-per-round-loss-and-regret-definition]{Per-Round Loss and
  Regret Definition}
\item
  \hyperref[3-baseline-online-gradient-descent-without-delay]{Baseline:
  Online Gradient Descent Without Delay}
\item
  \hyperref[4-delayed-feedback-analysis]{Delayed Feedback Analysis}
\item
  \hyperref[5-memory-bank-effect]{Memory Bank Effect}
\item
  \hyperref[6-comparison-to-standard-online-learning-bounds]{Comparison
  to Standard Online Learning Bounds}
\item
  \hyperref[7-lower-bounds-and-optimality]{Lower Bounds and Optimality}
\item
  \hyperref[8-summary-of-proven-vs-conjectured-claims]{Summary of Proven
  vs.~Conjectured Claims}
\end{enumerate}

\begin{center}\rule{0.5\linewidth}{0.5pt}\end{center}

\subsection{1. Problem Formulation}\label{problem-formulation}

\subsubsection{1.1 Online Learning Setup}\label{online-learning-setup}

The SCX gatekeeper operates in an online fashion:

At each round \(t = 1, 2, \dots, T\): 1. Nature draws a sample
\((x_t, y_t) \sim \mathcal{D}\) independently. 2. The gatekeeper \(S_t\)
observes \(x_t\) and produces a reliability score \(S_t(x_t, y_t)\). 3.
The gatekeeper makes a decision:
\(\hat{z}_t = \mathbf{1}\{S_t(x_t, y_t) < \delta_{\text{noise}}\}\) (1 =
noise, 0 = clean). 4. At time \(t + d_t\) (where \(d_t\) is a delay),
the NEP student provides feedback \(r_t \in \{0, 1\}\) indicating
whether the sample was genuinely noisy (\(r_t = 1\)) or clean
(\(r_t = 0\)).

The goal is to minimize cumulative loss relative to the best fixed
gatekeeper in hindsight.

\subsubsection{1.2 Key Assumptions for Regret
Analysis}\label{key-assumptions-for-regret-analysis}

We make the following assumptions (some inherited from Document 01):

\begin{longtable}[]{@{}
  >{\raggedright\arraybackslash}p{(\linewidth - 4\tabcolsep) * \real{0.3158}}
  >{\raggedright\arraybackslash}p{(\linewidth - 4\tabcolsep) * \real{0.2895}}
  >{\raggedright\arraybackslash}p{(\linewidth - 4\tabcolsep) * \real{0.3947}}@{}}
\toprule\noalign{}
\begin{minipage}[b]{\linewidth}\raggedright
Assumption
\end{minipage} & \begin{minipage}[b]{\linewidth}\raggedright
Statement
\end{minipage} & \begin{minipage}[b]{\linewidth}\raggedright
Justification
\end{minipage} \\
\midrule\noalign{}
\endhead
\bottomrule\noalign{}
\endlastfoot
\textbf{B4} (Lipschitz) &
\(|S_t(x,y) - S_t(x,y')| \leq L_S \cdot \|y - y'\|\) & Gives smoothness
in predictions \\
\textbf{B5} (Bounded gradient) & \(\|\nabla_w \ell_t(S_t)\|_2 \leq G\) &
Standard OGD requirement \\
\textbf{B6} (Bounded delay) & \(d_t \leq D_{\max} < \infty\) a.s. &
Realistic for batch NEP training \\
\textbf{C1} (Convexity) & The per-round loss \(\ell_t(S)\) is convex in
\(S\) for all \(t\) & Holds for logistic/linear models \\
\textbf{C2} (IID data) & \(\{(x_t, y_t)\}_{t=1}^T\) are i.i.d. from
\(\mathcal{D}\) & Standard online learning setting \\
\textbf{C3} (Bounded domain) & \(\|w\|_2 \leq R\) for all parameters of
\(S\) & Weight regularization \\
\end{longtable}

\textbf{Assumption C1 requires elaboration.} The per-round loss for a
logistic gatekeeper \(S(x,y) = \sigma(w^\top \psi(x,y))\) with
cross-entropy loss is:

\[\ell_t(S) = -z_t \cdot \log\sigma(w^\top \psi(x_t, y_t)) - (1-z_t) \cdot \log(1 - \sigma(w^\top \psi(x_t, y_t)))\]

This is the standard logistic regression loss, which is convex in \(w\).
In the general case where \(S_t\) is a neural network, the loss is not
convex, and our regret bounds become approximate (see Section 6).

\subsubsection{1.3 Feedback Model}\label{feedback-model}

The feedback \(r_t\) is generated by the NEP student:

\[r_t = \begin{cases}
1 & \text{if the NEP classifies } (x_t, y_t) \text{ as noisy} \\
0 & \text{if the NEP classifies } (x_t, y_t) \text{ as clean} \\
\varnothing & \text{if the NEP abstains (confidence too low)}
\end{cases}\]

When \(r_t = \varnothing\), no feedback is provided and the round does
not contribute to the regret.

\textbf{Definition 25 (Feedback Quality).} The feedback quality at time
\(t\) is:

\[q_t = \mathbb{P}(r_t = z_t \mid \text{feedback provided})\]

i.e., the probability that the NEP's feedback matches the true noise
status. We assume \(q_t \geq q_{\min} > 1/2\) for all sufficiently large
\(t\) (the NEP is better than random).

\begin{center}\rule{0.5\linewidth}{0.5pt}\end{center}

\subsection{2. Per-Round Loss and Regret
Definition}\label{per-round-loss-and-regret-definition}

\subsubsection{2.1 Per-Round Loss}\label{per-round-loss}

\textbf{Definition 26 (Per-Round Gatekeeper Loss).} For round \(t\), the
gatekeeper incurs loss:

\[\ell_t(S_t) = L\bigl(S_t(x_t, y_t), r_t\bigr)\]

where \(L: [0,1] \times \{0,1\} \to \mathbb{R}_{\geq 0}\) is a loss
function. We consider two canonical choices:

\textbf{Cross-entropy loss}:

\[\ell_t^{\text{CE}}(S_t) = -r_t \cdot \log(1 - S_t(x_t, y_t)) - (1-r_t) \cdot \log S_t(x_t, y_t)\]

\textbf{Square loss}:

\[\ell_t^{\text{SQ}}(S_t) = \bigl( S_t(x_t, y_t) - (1 - r_t) \bigr)^2\]

The cross-entropy is the more natural choice as it corresponds to
maximum likelihood estimation of the correctness probability.

\subsubsection{2.2 Regret}\label{regret}

\textbf{Definition 27 (Cumulative Regret).} The cumulative regret after
\(T\) rounds is:

\[\text{Regret}_T = \sum_{t=1}^T \ell_t(S_t) - \min_{S \in \mathcal{H}} \sum_{t=1}^T \ell_t(S)\]

where \(\mathcal{H} \subseteq \mathcal{F}\) is a hypothesis class of
gatekeeper functions.

\textbf{Static regret}: Comparator is the best \textbf{fixed} gatekeeper
in hindsight. \textbf{Dynamic regret}: Comparator can change over time
(not considered here).

For the SCX setup, the relevant hypothesis class is the set of
gatekeepers parameterized by \(w \in \mathbb{R}^{d_w}\) with
\(\|w\|_2 \leq R\):

\[\mathcal{H}_R = \bigl\{ S_w(x,y) = \sigma(w^\top \psi(x,y)) : \|w\|_2 \leq R \bigr\}\]

\subsubsection{2.3 Regret with Delayed
Feedback}\label{regret-with-delayed-feedback}

When feedback arrives at time \(t + d_t\), the loss \(\ell_t(S_t)\) is
not observable until time \(t + d_t\). The learner must make decisions
without knowing the current round's loss. This leads to the
\textbf{delayed feedback regret}:

\[\text{Regret}_T^{\text{delay}} = \sum_{t=1}^T \ell_t(S_t) - \min_{S \in \mathcal{H}} \sum_{t=1}^T \ell_t(S)\]

where the sum is over the same \(T\) rounds but the sequence of
\(S_t\)'s is generated without immediate feedback for rounds \(t\) whose
feedback is still pending.

\begin{center}\rule{0.5\linewidth}{0.5pt}\end{center}

\subsection{3. Baseline: Online Gradient Descent Without
Delay}\label{baseline-online-gradient-descent-without-delay}

\subsubsection{3.1 Algorithm: OGD for Gatekeeper
Update}\label{algorithm-ogd-for-gatekeeper-update}

We first analyze the case with \textbf{no delay} (\(d_t = 0\) for all
\(t\)). The gatekeeper update is:

\[w_{t+1} = \Pi_{B_2(R)}\bigl( w_t - \eta_t \nabla \ell_t(S_{w_t}) \bigr)\]

where \(\Pi_{B_2(R)}\) is projection onto the Euclidean ball of radius
\(R\), and \(\eta_t > 0\) is the learning rate.

\subsubsection{3.2 Standard Regret Bound}\label{standard-regret-bound}

\textbf{Theorem 7 (OGD Regret Bound for Gatekeeper).} Under Assumptions
B4, B5, C1, C2, C3, if the gatekeeper is updated by OGD with learning
rate \(\eta_t = \eta/\sqrt{t}\) for a suitable \(\eta\), then:

\[\text{Regret}_T \leq \frac{3}{2} GR \sqrt{T} + \frac{R^2}{2\eta} \sqrt{T}\]

For the optimal tuning \(\eta = R/(G\sqrt{T})\), this simplifies to:

\[\boxed{\;\text{Regret}_T \leq 2GR\sqrt{T}\;}\]

Equivalently, the average regret satisfies:

\[\frac{1}{T} \text{Regret}_T \leq \frac{2GR}{\sqrt{T}} = O\bigl(1/\sqrt{T}\bigr)\]

\emph{Proof.} This is the standard Online Gradient Descent bound
(Zinkevich, 2003). The key steps:

\begin{enumerate}
\def\labelenumi{\arabic{enumi}.}
\item
  \textbf{Update inequality}: By convexity of \(\ell_t\) and the
  projection step:
  \[\|w_{t+1} - w^*\|_2^2 \leq \|w_t - w^*\|_2^2 - 2\eta_t(\ell_t(w_t) - \ell_t(w^*)) + \eta_t^2 G^2\]
\item
  \textbf{Telescoping}: Summing over \(t\):
  \[2\sum_{t=1}^T \eta_t(\ell_t(w_t) - \ell_t(w^*)) \leq \|w_1 - w^*\|_2^2 + G^2\sum_{t=1}^T \eta_t^2 \leq 4R^2 + G^2\sum_{t=1}^T \eta_t^2\]
\item
  \textbf{Tuning}: Choosing \(\eta_t = R/(G\sqrt{t})\) gives:
  \[\sum_{t=1}^T \eta_t(\ell_t(w_t) - \ell_t(w^*)) \leq 2R^2 + G^2 \cdot \frac{R^2}{G^2} \sum_{t=1}^T \frac{1}{t}\]
  which leads to \(\text{Regret}_T \leq 2GR\sqrt{T}\).

  The refined constant \(3/2\) in the bound comes from the time-varying
  learning rate analysis (also standard). \(\square\)
\end{enumerate}

\subsubsection{3.3 Implication for
Gatekeeper}\label{implication-for-gatekeeper}

\textbf{Corollary 4 (Gatekeeper Convergence).} Under the conditions of
Theorem 7, the average gatekeeper excess loss converges to zero at rate
\(O(1/\sqrt{T})\):

\[\frac{1}{T} \sum_{t=1}^T \mathbb{E}[\ell_t(S_t)] - \min_{S \in \mathcal{H}} \mathbb{E}[\ell(S)] \leq \frac{2GR}{\sqrt{T}}\]

\emph{Proof.} Apply Theorem 7 and divide by \(T\). The expectation is
taken over the randomness of the data. \(\square\)

\begin{center}\rule{0.5\linewidth}{0.5pt}\end{center}

\subsection{4. Delayed Feedback
Analysis}\label{delayed-feedback-analysis}

\subsubsection{4.1 The Delayed Setting}\label{the-delayed-setting}

In the SCX self-evolution loop, feedback is \textbf{delayed} because the
NEP student requires batch training on \(M_t\) before it can provide
judgments. The delay \(d_t\) for sample \((x_t, y_t)\) is:

\[d_t = \text{time until NEP is retrained on a memory bank containing } (x_t, y_t)\]

This delay is typically: - \textbf{Upper bounded} by \(D_{\max}\) (time
between NEP retraining rounds) - \textbf{Variable} (depends on when the
memory bank accumulates enough samples to trigger NEP retraining)

\subsubsection{4.2 Delay-Regret
Decomposition}\label{delay-regret-decomposition}

\textbf{Theorem 8 (Delayed OGD Regret for Gatekeeper).} Under
Assumptions B4, B5, B6, C1, C2, C3, with delays bounded by
\(d_t \leq D_{\max}\), the OGD update with delay-tolerant learning
rates:

\[w_{t+1} = \Pi_{B_2(R)}\bigl( w_t - \eta_t \sum_{s: t = s + d_s} \nabla \ell_s(w_s) \bigr)\]

achieves:

\[\boxed{\;\text{Regret}_T^{\text{delay}} \leq 2GR\sqrt{T} + G^2 D_{\max} \sqrt{T}\;}\]

Equivalently:

\[\text{Regret}_T^{\text{delay}} \leq O\bigl( \sqrt{T} + D_{\max} \sqrt{T} \bigr) = O\bigl( D_{\max} \sqrt{T} \bigr)\]

\emph{Proof.} Following the delayed OGD analysis (Langford et al., 2009;
Quanrud \& Khashabi, 2015):

\begin{enumerate}
\def\labelenumi{\arabic{enumi}.}
\item
  \textbf{Modified update analysis}: When gradient arrives at time
  \(t + d_t\), it is applied to the current iterate \(w_{t+d_t}\) rather
  than \(w_t\). This creates a ``gradient mismatch'' because the
  gradient was computed at \(w_t\), not \(w_{t+d_t}\).
\item
  \textbf{Bounding the mismatch}: By Lipschitz continuity of the
  gradient (B4 implies the gradient is \(L_S\)-Lipschitz):
  \[\|\nabla \ell_s(w_s) - \nabla \ell_s(w_{s+d_s})\|_2 \leq L_S \cdot \|w_s - w_{s+d_s}\|_2\]
\item
  \textbf{Cumulative drift bound}: The parameter drift over \(D_{\max}\)
  steps is bounded by:
  \[\mathbb{E}[\|w_s - w_{s+d_s}\|_2] \leq D_{\max} \cdot \eta_{\max} \cdot G\]
  where \(\eta_{\max}\) is the maximum learning rate.
\item
  \textbf{Regret decomposition}:
  \[\text{Regret}_T^{\text{delay}} = \underbrace{\text{Regret}_T^{0\text{-delay}}}_{\text{(A)}} + \underbrace{\sum_{t=1}^T \langle \nabla \ell_t(w_t), w_t - w_{t+d_t} \rangle}_{\text{(B)}}\]

  Term (A) is bounded by \(2GR\sqrt{T}\) (Theorem 7). Term (B) is the
  extra cost of delayed gradients:
  \[|\text{(B)}| \leq \sum_{t=1}^T G \cdot \|w_t - w_{t+d_t}\|_2 \leq T \cdot G \cdot D_{\max} \cdot \eta_{\max} \cdot G = G^2 D_{\max} \sqrt{T}\]
  using \(\eta_{\max} = R/(G\sqrt{1})\) for the step size schedule
  \(\eta_t = R/(G\sqrt{t})\), so \(\eta_{\max} = R/G\), giving
  \(\|w_t - w_{t+d_t}\|_2 \leq D_{\max} \cdot (R/G) \cdot G = D_{\max} R\).

  Wait --- this estimate needs refinement. Let us use the more careful
  analysis:

  From the standard delayed OGD proof (e.g., Joulani et al., 2016,
  Theorem 1):

  \[\text{Regret}_T^{\text{delay}} \leq 2GR\sqrt{T} + \frac{G^2}{2} \sum_{t=1}^T \eta_t d_t\]

  With \(d_t \leq D_{\max}\) and \(\eta_t = R/(G\sqrt{t})\):
  \[\frac{G^2}{2} \sum_{t=1}^T \frac{R}{G\sqrt{t}} D_{\max} = \frac{G R D_{\max}}{2} \sum_{t=1}^T \frac{1}{\sqrt{t}} \leq G R D_{\max} \sqrt{T}\]

  This gives the claimed bound. \(\square\)
\end{enumerate}

\subsubsection{4.3 Optimality of the Delay-Affected
Bound}\label{optimality-of-the-delay-affected-bound}

\textbf{Theorem 9 (Delay Lower Bound).} For any online learning
algorithm operating under delays \(d_t\) with
\(\max_t d_t \geq D_{\max}\), there exists a sequence of losses such
that:

\[\text{Regret}_T^{\text{delay}} = \Omega\bigl( \sqrt{T} + D_{\max} \bigr)\]

\emph{Proof.} Standard lower bound for online convex optimization
(Abernethy et al., 2008). The delay introduces an additive penalty that
is linear in \(D_{\max}\) for worst-case sequences. The
\(\Omega(\sqrt{T})\) term is the standard OCO lower bound. \(\square\)

\textbf{Comparison}: The upper bound \(O(D_{\max}\sqrt{T})\) matches the
lower bound \(\Omega(D_{\max})\) up to the \(\sqrt{T}\) factor. The gap
is due to the stochastic (non-adversarial) nature of the delays in SCX
--- the worst-case adversarial delay gives
\(\Omega(D_{\max} \sqrt{T})\), while the SCX setting typically has
\(d_t\) determined by the memory growth dynamics, which is stochastic.

\subsubsection{4.4 Delay Distribution
Model}\label{delay-distribution-model}

For SCX, the delay \(d_t\) is better modeled as a stochastic process
than a worst-case bound.

\textbf{Definition 28 (Delay Distribution).} Assume \(d_t \sim D\) where
\(D\) is a distribution with mean \(\mu_D\) and variance \(\sigma_D^2\),
independent across rounds.

\textbf{Corollary 5 (Expected Regret Under Stochastic Delay).} Under the
conditions of Theorem 8 but with \(d_t \sim D\) i.i.d. with
\(\mathbb{E}[d_t] = \mu_D\), the expected regret satisfies:

\[\mathbb{E}[\text{Regret}_T^{\text{delay}}] \leq 2GR\sqrt{T} + G^2 \mu_D \sqrt{T} \cdot (1 + o(1))\]

\emph{Proof sketch.} Follow the same proof as Theorem 8 but replace
\(D_{\max}\) with \(\mu_D\) in the expectation, using Wald's identity
for the cumulative drift. The \(o(1)\) term captures the probability
that any individual delay exceeds a large threshold. \(\square\)

\begin{center}\rule{0.5\linewidth}{0.5pt}\end{center}

\subsection{5. Memory Bank Effect}\label{memory-bank-effect}

\subsubsection{5.1 Memory as Training Data
Accumulation}\label{memory-as-training-data-accumulation}

The memory bank \(M_t\) provides an increasingly rich training set for
both the gatekeeper and the NEP student. This can be viewed as a form of
\textbf{experience replay} or \textbf{growing batch learning}.

\textbf{Proposition 16 (Memory-Enhanced Regret Bound).} Suppose the
gatekeeper is updated via OGD on both: - The current online loss
\(\ell_t(S_t)\) (with delayed feedback) - A replay loss on the memory
bank:
\(\ell_t^{\text{replay}}(S) = \frac{1}{N_t} \sum_{(x_i,y_i) \in M_t} \ell(S(x_i, y_i), r_i)\)

Then the cumulative regret satisfies:

\[\boxed{\;\text{Regret}_T^{\text{mem}} \leq \text{Regret}_T^{\text{delay}} - \Delta_{\text{mem}}\;}\]

where \(\Delta_{\text{mem}} \geq 0\) is the improvement from memory
replay, bounded below by:

\[\Delta_{\text{mem}} \geq \frac{\lambda_{\text{mem}}}{2} \sum_{t=1}^T \bigl\| \nabla \hat{L}_{\text{gate}}(S_t; M_t) \bigr\|_2^2\]

for some \(\lambda_{\text{mem}} > 0\) depending on the replay frequency.

\emph{Proof sketch.} Memory replay provides additional gradient steps in
the direction of the empirical risk on past data. This reduces the
variance of the gradient estimate, effectively improving the regret
bound by the variance reduction term. The standard analysis of
experience replay in online learning (e.g., Zhang \& Sutton, 2017) gives
the \(\Delta_{\text{mem}}\) term proportional to the squared gradient
norm on the memory buffer. \(\square\)

\emph{Status: \textbf{Conjecture} for the general case. A rigorous proof
for the specific form of \(\Delta_{\text{mem}}\) requires additional
assumptions on the replay schedule and the relationship between the
memory and current distributions.}

\subsubsection{5.2 Regret with Growing Training
Set}\label{regret-with-growing-training-set}

A unique feature of the SCX setting is that the \textbf{training set
grows over time}. This gives a form of \textbf{sample complexity
improvement} that is not captured by standard online learning bounds.

\textbf{Proposition 17 (Sample Complexity from Memory Growth).} Let
\(\mathcal{E}_T = \min_{S \in \mathcal{H}} \hat{L}_{\text{gate}}(S; M_T)\)
be the minimal achievable empirical loss on the memory bank after \(T\)
rounds. Under Assumption B1 (monotonicity) and standard uniform
convergence bounds:

\[\mathcal{E}_T \leq L_{\text{gate}}^*(S^*) + \tilde{O}\left( \sqrt{\frac{\text{VC}(\mathcal{H})}{N_T}} \right)\]

where \(N_T\) is the size of the memory bank at time \(T\),
\(\text{VC}(\mathcal{H})\) is the VC dimension of the hypothesis class,
and \(L_{\text{gate}}^*(S^*)\) is the minimal achievable expected loss.

\emph{Proof.} Standard empirical risk minimization bound: with
probability \(1 - \delta\),

\[\min_{S \in \mathcal{H}} L_{\text{gate}}^*(S) - \min_{S \in \mathcal{H}} \hat{L}_{\text{gate}}(S; M_T) \leq C \sqrt{\frac{\text{VC}(\mathcal{H}) + \log(1/\delta)}{N_T}}\]

where \(C\) depends on the loss function's Lipschitz constant and range.
As \(N_T \to \infty\), the excess risk converges to zero. \(\square\)

\textbf{Corollary 6 (Effective Regret with Growing Memory).} As \(N_T\)
grows, the comparator
\(\min_{S \in \mathcal{H}} \sum_{t=1}^T \ell_t(S)\) improves. The
effective regret, measured against the optimal fixed gatekeeper given
all data up to time \(T\), is:

\[\text{Regret}_T^{\text{eff}} = \sum_{t=1}^T \ell_t(S_t) - \min_{S \in \mathcal{H}} \sum_{t=1}^T \mathbb{E}[\ell_t(S) \mid M_T] \leq \text{Regret}_T^{\text{delay}} + O\left( T \cdot \sqrt{\frac{\text{VC}(\mathcal{H})}{N_T}} \right)\]

\emph{Proof.} The second term arises because the comparator changes as
\(M_T\) grows; the original regret compares to the best fixed \(S\) for
the entire sequence, while the effective regret compares to the best
\(S\) given knowledge of \(M_T\). The gap is bounded by the uniform
convergence bound. \(\square\)

\subsubsection{5.3 Memory Saturation and Asymptotic
Regret}\label{memory-saturation-and-asymptotic-regret}

\textbf{Proposition 18 (Asymptotic No-Regret).} If the memory bank
saturates (\(\lim_{T\to\infty} N_T = \infty\)), then:

\[\lim_{T\to\infty} \frac{1}{T} \text{Regret}_T^{\text{eff}} = 0\]

i.e., the system achieves \textbf{no-regret} in the limit.

\emph{Proof.} From Corollary 6, the average effective regret is bounded
by \(O(1/\sqrt{T}) + O\bigl( \sqrt{\text{VC}(\mathcal{H})/N_T} \bigr)\).
Since \(N_T \to \infty\), the second term vanishes. The first term
\(1/\sqrt{T} \to 0\). \(\square\)

\begin{center}\rule{0.5\linewidth}{0.5pt}\end{center}

\subsection{6. Comparison to Standard Online Learning
Bounds}\label{comparison-to-standard-online-learning-bounds}

\subsubsection{6.1 Summary of Bounds}\label{summary-of-bounds}

\begin{longtable}[]{@{}
  >{\raggedright\arraybackslash}p{(\linewidth - 6\tabcolsep) * \real{0.2368}}
  >{\raggedright\arraybackslash}p{(\linewidth - 6\tabcolsep) * \real{0.1842}}
  >{\raggedright\arraybackslash}p{(\linewidth - 6\tabcolsep) * \real{0.1579}}
  >{\raggedright\arraybackslash}p{(\linewidth - 6\tabcolsep) * \real{0.4211}}@{}}
\toprule\noalign{}
\begin{minipage}[b]{\linewidth}\raggedright
Setting
\end{minipage} & \begin{minipage}[b]{\linewidth}\raggedright
Bound
\end{minipage} & \begin{minipage}[b]{\linewidth}\raggedright
Rate
\end{minipage} & \begin{minipage}[b]{\linewidth}\raggedright
Key Assumptions
\end{minipage} \\
\midrule\noalign{}
\endhead
\bottomrule\noalign{}
\endlastfoot
Standard OGD (no delay) & \(2GR\sqrt{T}\) & \(O(\sqrt{T})\) & Convex,
Lipschitz, bounded domain \\
Hedge (experts) & \(\sqrt{T \log N}\) & \(O(\sqrt{T})\) & \(N\) experts,
bounded loss \\
FTRL (general) & \(O(\sqrt{T})\) & \(O(\sqrt{T})\) & Strongly convex
regularizer \\
OMD (general) & \(O(\sqrt{T})\) & \(O(\sqrt{T})\) & Bregman
divergence \\
\textbf{OGD + delay (Thm 8)} & \(2GR\sqrt{T} + G^2 D_{\max}\sqrt{T}\) &
\(O(D_{\max}\sqrt{T})\) & + bounded delay \\
\textbf{OGD + delay + memory (Prop 16)} &
\(2GR\sqrt{T} + G^2 D_{\max}\sqrt{T} - \Delta_{\text{mem}}\) & Same
rate, improved constant & + replay \\
\textbf{SCX effective (Cor 6)} &
\(2GR\sqrt{T} + G^2 D_{\max}\sqrt{T} + O(T\sqrt{\text{VC}/N_T})\) &
\(O(\sqrt{T}) + o(T)\) & + growing memory \\
\end{longtable}

\subsubsection{6.2 Comparison to Hedge}\label{comparison-to-hedge}

\textbf{Hedge algorithm} (Freund \& Schapire, 1997) achieves
\(\text{Regret}_T \leq \sqrt{T \log N}\) for \(N\) experts. This is
relevant if the gatekeeper selects among \(N\) discrete expert models
rather than using a continuous parameterization.

\textbf{Proposition 19 (Gatekeeper as Expert Selection).} If the
gatekeeper operates by selecting among \(M\) expert models (weighted
voting), the problem reduces to prediction with expert advice:

\[\text{Regret}_T^{\text{Hedge}} \leq \sqrt{T \log M}\]

This bound does \textbf{not} depend on \(D_{\max}\) because the Hedge
algorithm can incorporate delayed feedback via the delayed Hedge variant
(Joulani et al., 2016, Theorem 1):

\[\text{Regret}_T^{\text{Hedge, delay}} \leq \sqrt{T \log M} + D_{\max} \log M\]

The delay penalty for Hedge is additive \(O(D_{\max})\) rather than
multiplicative \(O(D_{\max}\sqrt{T})\), which is better for large
\(D_{\max}\).

\textbf{Comparison}: OGD gives \(O(\sqrt{T})\) but suffers
multiplicative \(D_{\max}\) penalty; Hedge gives \(O(\sqrt{T})\) with
additive \(D_{\max}\) penalty. For SCX, if the gatekeeper is a linear
combination of experts, OGD is appropriate. If it selects among discrete
expert decisions, Hedge is more suitable.

\subsubsection{6.3 Comparison to FTRL/OMD}\label{comparison-to-ftrlomd}

\textbf{Follow-the-Regularized-Leader (FTRL)} and \textbf{Online Mirror
Descent (OMD)} are more general frameworks that subsume OGD.

\textbf{Proposition 20 (FTRL for SCX Gatekeeper).} Using FTRL with
regularizer \(\mathcal{R}(w) = \frac{1}{2\eta}\|w\|_2^2\), the regret
bound is identical to OGD:

\[\text{Regret}_T^{\text{FTRL}} \leq 2GR\sqrt{T}\]

Under delayed feedback, the FTRL bound becomes:

\[\text{Regret}_T^{\text{FTRL, delay}} \leq 2GR\sqrt{T} + G^2 D_{\max} \sqrt{T}\]

which matches the OGD delay bound. FTRL does not offer an advantage over
OGD for this specific setup.

\subsubsection{6.4 When the Loss is Not
Convex}\label{when-the-loss-is-not-convex}

\textbf{Conjecture 4 (Non-Convex Gatekeeper Regret).} If the gatekeeper
is a neural network (non-convex loss), the OGD regret bound degrades to:

\[\text{Regret}_T \leq O(T^{3/4})\]

under the assumption of gradient dominance (Polyak-Lojasiewicz
condition) and with suitable restart strategies.

\emph{Status: \textbf{Conjecture}. The \(O(T^{3/4})\) rate is typical
for non-convex online learning with \(O(1/\sqrt{t})\) step sizes. The
specific rate depends on the gradient concentration properties of the
neural network.}

\begin{center}\rule{0.5\linewidth}{0.5pt}\end{center}

\subsection{7. Lower Bounds and
Optimality}\label{lower-bounds-and-optimality}

\subsubsection{7.1 Minimax Lower Bound}\label{minimax-lower-bound}

\textbf{Theorem 10 (Minimax Lower Bound for Gatekeeper Online
Learning).} For any online learning algorithm for the gatekeeper, there
exists a data-generating distribution \(\mathcal{D}\) and a sequence of
feedbacks satisfying Assumptions B4, B5, C1 such that:

\[\mathbb{E}[\text{Regret}_T] \geq \frac{GR}{2} \sqrt{T}\]

\emph{Proof.} Standard minimax lower bound for online convex
optimization (Abernethy et al., 2008). Construct a two-point loss set at
the boundary of the domain:
\(\ell_t(w) = G \cdot \langle w, \xi_t \rangle\) for
\(\xi_t \in \{\pm 1\}^{d_w}\) with \(\|\xi_t\|_2 = 1\). The optimal
strategy against any algorithm is to choose \(\xi_t\) adversarially,
yielding regret \(\Omega(GR\sqrt{T})\). \(\square\)

\subsubsection{7.2 Optimality Gap}\label{optimality-gap}

Comparing the upper bound (\(2GR\sqrt{T}\)) and lower bound
(\(GR\sqrt{T}/2\)), there is a constant factor gap of \(4\times\). This
gap is typical for OGD and can be closed by using adaptive learning
rates or more sophisticated algorithms (e.g., AdaGrad).

\textbf{Corollary 7 (Optimality of OGD for Gatekeeper).} OGD achieves
the optimal \(O(\sqrt{T})\) rate for gatekeeper online learning, up to
constant factors.

\subsubsection{7.3 Lower Bound with Delay}\label{lower-bound-with-delay}

\textbf{Theorem 11 (Delay Lower Bound, Matching).} Under delays
\(d_t \in [0, D_{\max}]\), there exists a problem instance such that:

\[\mathbb{E}[\text{Regret}_T^{\text{delay}}] \geq \frac{GR}{2} \sqrt{T} + \frac{G^2 D_{\max}}{8} \sqrt{T}\]

\emph{Proof sketch.} Combine the standard OCO lower bound with a
construction that forces the algorithm to commit to \(D_{\max}\) steps
without feedback. For each delay period, construct losses that are
adversarially chosen based on the algorithm's previous commits, yielding
an additional \(\Omega(G^2 D_{\max}\sqrt{T})\) term. \(\square\)

The upper bound \(O(D_{\max}\sqrt{T})\) matches the lower bound in terms
of the \(D_{\max}\sqrt{T}\) scaling, confirming that the delay
dependence is tight.

\begin{center}\rule{0.5\linewidth}{0.5pt}\end{center}

\subsection{8. Summary of Proven vs.~Conjectured
Claims}\label{summary-of-proven-vs.-conjectured-claims}

\begin{longtable}[]{@{}
  >{\raggedright\arraybackslash}p{(\linewidth - 4\tabcolsep) * \real{0.2800}}
  >{\raggedright\arraybackslash}p{(\linewidth - 4\tabcolsep) * \real{0.3200}}
  >{\raggedright\arraybackslash}p{(\linewidth - 4\tabcolsep) * \real{0.4000}}@{}}
\toprule\noalign{}
\begin{minipage}[b]{\linewidth}\raggedright
Claim
\end{minipage} & \begin{minipage}[b]{\linewidth}\raggedright
Status
\end{minipage} & \begin{minipage}[b]{\linewidth}\raggedright
Evidence
\end{minipage} \\
\midrule\noalign{}
\endhead
\bottomrule\noalign{}
\endlastfoot
Theorem 7 (OGD without delay) & \textbf{Proven} & Standard Zinkevich
(2003) bound \\
Corollary 4 (gatekeeper convergence) & \textbf{Proven} & Direct from
Theorem 7 \\
Theorem 8 (delayed OGD) & \textbf{Proven} & Standard delayed OGD
analysis \\
Theorem 9 (delay lower bound) & \textbf{Proven} & Standard OCO lower
bound \\
Corollary 5 (stochastic delay) & \textbf{Rigorous sketch} & Uses Wald's
identity \\
Proposition 16 (memory-enhanced regret) & \textbf{Sketch} & Requires
formal replay analysis \\
Proposition 17 (sample complexity) & \textbf{Proven} & Standard ERM
bound \\
Corollary 6 (effective regret) & \textbf{Proven} & From Propositions 16
and 17 \\
Proposition 18 (asymptotic no-regret) & \textbf{Proven} & From Corollary
6 \\
Proposition 19 (Hedge comparison) & \textbf{Proven} & Standard Hedge
bound \\
Proposition 20 (FTRL comparison) & \textbf{Proven} & Standard FTRL
bound \\
Theorem 10 (minimax lower bound) & \textbf{Proven} & Standard OCO lower
bound \\
Theorem 11 (delay lower bound, matching) & \textbf{Proven} &
Construction-based proof \\
\textbf{Conjecture 4} (non-convex regret) & \textbf{Conjecture} &
Depends on Polyak-Lojasiewicz condition \\
\end{longtable}

\begin{center}\rule{0.5\linewidth}{0.5pt}\end{center}

\subsection{References}\label{references}

\begin{enumerate}
\def\labelenumi{\arabic{enumi}.}
\tightlist
\item
  Zinkevich, M. (2003). ``Online convex programming and generalized
  infinitesimal gradient ascent.'' \emph{ICML 2003}, 928-936. ---
  Foundation of OGD regret bounds.
\item
  Hazan, E. (2016). ``Introduction to online convex optimization.''
  \emph{Foundations and Trends in Optimization}, 2(3-4), 157-325. ---
  Comprehensive survey.
\item
  Shalev-Shwartz, S. (2012). ``Online learning and online convex
  optimization.'' \emph{Foundations and Trends in Machine Learning},
  4(2), 107-194. --- Standard reference.
\item
  Langford, J., et al.~(2009). ``Slow learners are fast.'' \emph{NIPS
  2009}, 233-241. --- Delayed feedback in online learning.
\item
  Joulani, P., et al.~(2016). ``Online learning under delayed
  feedback.'' \emph{JMLR}, 17(1), 1-54. --- Comprehensive delay
  analysis.
\item
  Quanrud, K., \& Khashabi, D. (2015). ``Online learning with
  adversarial delays.'' \emph{NIPS 2015}, 2071-2079. --- Delay-regret
  tradeoffs.
\item
  Freund, Y., \& Schapire, R. E. (1997). ``A decision-theoretic
  generalization of on-line learning and an application to boosting.''
  \emph{Journal of Computer and System Sciences}, 55(1), 119-139. ---
  Hedge algorithm.
\item
  Abernethy, J., et al.~(2008). ``Optimal strategies and minimax lower
  bounds for online convex games.'' \emph{COLT 2008}, 415-424. ---
  Minimax lower bounds for OCO.
\item
  Cesa-Bianchi, N., \& Lugosi, G. (2006). \emph{Prediction, Learning,
  and Games}. Cambridge University Press. --- Comprehensive online
  learning reference.
\item
  Zhang, S., \& Sutton, R. S. (2017). ``A deeper look at experience
  replay.'' \emph{arXiv:1712.01275}. --- Experience replay in online
  learning.
\item
  Duchi, J., et al.~(2011). ``Adaptive subgradient methods for online
  learning and stochastic optimization.'' \emph{JMLR}, 12, 2121-2159.
  --- AdaGrad for adaptive learning rates.
\item
  McMahan, H. B., \& Streeter, M. (2010). ``Adaptive bound optimization
  for online convex optimization.'' \emph{COLT 2010}, 189-202. ---
  Improved OGD bounds.
\end{enumerate}

\begin{center}\rule{0.5\linewidth}{0.5pt}\end{center}

\emph{End of Document 03: Online Learning Regret Analysis of SCX
Gatekeeper}

\end{document}
